% !TEX root = master.tex
% !TEX program = lualatex
\lecture{5}{19-03-2026}{Singularities}{}

\begin{remark}
The Laurent series of a function is unique. Hence, if we compute a power series expansion for $f$ (e.g. from a known expansion), we have determined its Laurent series.
\end{remark}

\begin{parag}{Example}
  

We know:
\[
\frac{1}{1-z} = \sum_{n=0}^{\infty} z^n, \quad |z|<1
\]

Then:
\[
\frac{1}{1+z^2} = \frac{1}{1-(-z^2)} = \sum_{n=0}^{\infty} (-1)^n z^{2n}
\]

So:
\[
\frac{1}{1+z^2} = \sum_{k=0}^{\infty} c_k z^k
\]
where
\[
c_k =
\begin{cases}
0 & \text{if $k$ odd} \\
(-1)^{k/2} & \text{if $k$ even}
\end{cases}
\]

\end{parag}

\begin{parag}{Example: Laurent expansion around $z_0$}

Let $f: \mathbb{C} \setminus \{0, -1 \} \to \mathbb{C}$
\[
f(z) = \frac{1}{z(z+1)}
\]

We study the Laurent expansion at $z_0$.

\textbf{Case 1: $z_0 = 0$}

\[
f(z) = \frac{1}{z} - \frac{1}{z+1}
\]

We expand:
\[
\frac{1}{z+1} = \sum_{n=0}^{\infty} (-1)^n z^n, \quad |z|<1
\]

Thus:
\[
f(z) = \frac{1}{z} - \sum_{n=0}^{\infty} (-1)^n z^n
= \sum_{n=-1}^{\infty} (-1)^{n+1} z^n
\]

Radius of convergence:
\[
R = 1
\]
as the nearest singularity is at $z = -1$
(distance of $z_0$ to the next distinct singular point of $f$)
\textbf{Case 2: $z_0 = -1$}

\[
f(z) = -\frac{1}{z+1} + \frac{1}{z}
\]

Expand around $z=-1$:
\[
\frac{1}{z} = \frac{1}{(z+1)-1} = 
\frac{1}{-1(1-(z+1))}= \frac{-1}{1-(z+1)}=\sum_{n=0}^{\infty} (-1) (z+1)^n, \quad |z+1|<1
\]

So:
\[
f(z) = -\frac{1}{z - (-1)} + \sum_{n=0}^{\infty} (-1) (z+1)^n = \sum_{n=1}^{\infty}(-1)(z+1)^n 
\]

Radius of convergence:
\[
R = 1
\]
(Here the center = $z_0$ and the nearest singularity is at z = 0 hence the distance is $\left|0 - (-1)\right| = 1 $ ) 
\textbf{Case 3: $z_0 \neq 0, -1$}

$f$ is holomorphic near $z_0$, so $f$ can be written as a regular Taylor series. Converges in any disk around $z_0$ which does not contains neither 0 and -1. 

Radius:
\[
R = \min\{|z_0|, |z_0 + 1|\}
\]

\end{parag}

\section{Types of Singularities}

\subsection{Regular Point}

\begin{definition}
$f$ has a \textbf{regular point} at $z_0$ if the singular part of its Laurent series at $z_0$ vanishes.
\end{definition}

\begin{remark}
$f$ is holomorphic at $z_0$ $\iff$ $z_0$ is a regular point.

$z_0$ is a regular point $\iff lim_{z \to z_0} f(z)$ exists
\end{remark}

\begin{example}
\[
f(z) = e^z = \sum_{n=0}^{\infty} \frac{z^n}{n!}
\]
$z=0$ is a regular point (no terms of the form $z^{-1}, z^{-2}, \dots$ ($n < 0$))
\end{example}

\subsection{Pole}

\begin{definition}
$f$ has a \textbf{pole of order $m$} at $z_0$ if for $n < -m, c_n = 0, c_{-m}\neq 0$:
\[
f(z) = \frac{c_{-m}}{(z-z_0)^m} + \cdots + \frac{c_{-1}}{z-z_0} + \text{regular terms}
\]
\end{definition}

\begin{example}
\[
f(z) = \frac{1}{z} \quad \Rightarrow \text{pole of order 1 at } z=0
\]

\[
f(z) = \frac{1}{(z+2)^2} + \frac{1}{z}
\]
pole of order 2 at $z=-2$, pole of order 1 at $z=0$.
\end{example}

\begin{remark}
$f$ has a pole of order $m$ at $z_0$ $\iff$ $\lim_{z \to z_0} (z-z_0)^m f(z)$

exist and is not zero.
\end{remark}
In particular if $f$ has a pole of order $m$ at $z_0$ we can write \[
  f(z) = \frac{g(z)}{(z-z_0)^m}
\] 
where $g$ is holo. and $g(z_0) \neq 0$ (opposite is also true).

\subsection{Essential Singularity}

\begin{definition}
$f$ has an \textbf{essential singularity} at $z_0$ if its Laurent series has infinitely many coefficient $c_n \neq 0$ with $n \leq -1$.
\end{definition}

\begin{example}

$f(z) = e^{1/z}$ at $z_0 = 0$ 
\[
f(z) = \sum_{n=0}^{\infty} \frac{1}{n!} z^{-n} = \sum_{n= -\infty}^{0}\frac{1}{(-n)!}z^n 
\]
Essential singularity at $z=0$.
\end{example}

\section{Quotients of Taylor Series}

Let $p, q: D \setminus \{z_0\} \to \mathbb{C}$ holo with $z_0$ being regular for both. Write their Laurent series as:
\[
p(z) = a_0 + a_1(z-z_0) + \dots + a_k (z-z_0)^k
\quad
q(z) = b_0 + b_1(z-z_0) + \dots + b_\ell (z-z_0)^\ell
\]

Consider:
\[
f(z) = \frac{p(z)}{q(z)}
\]

if $b_0 \neq 0$ so $q(z_0) \neq 0$. We have \[
  \lim_{z \to z_0}\frac{p(z)}{q(z)} = \frac{p(z_0)}{q(z_0)} = \frac{a_0}{b_0} \quad \text{exists, so } z_0 \text{ is a regular point for } \frac{p}{q}
\] 

if $a_0 = b_0$ but $b_1 \neq 0$ then we have \[
  \lim_{z \to z_0} \frac{a_1(z - z_0) + a_2(z-z_0)^2 + \dots}{b_1(z-z_0) + b_2(z-z_0)^2 + \dots} = \lim_{z \to z_0} \frac{a_1 + a_2(z-z_0) + \dots}{b_1 + b_2(z-z_0) + \dots} = \frac{a_1}{b_1}
\]exists hence $z_0$ is a regular point for $f$. (similarly if $a_0 = b_0 = 0$ and $a_1 = b_1 = 0$ yet $b_2 \neq  0$) 

In summary, either expand in power series and take out as many common powers of $z-z_o$ from the nominator and denominator, and take the \textbf{limit} or us \textbf{l'Hôpital's rule}

\begin{itemize}
\item If $k \ge \ell$: regular point, as the limit = \begin{cases}
    0 & \text{if } k>\ell \\
    $\frac{a_k}{b_\ell}$ & \text{if } k=\ell
\end{cases} exists
\item If $k < \ell$: pole of order $\ell - k$
\end{itemize}

\paragraph{Example}
\[
f(z) = \frac{z}{\sin z}
\]where $z_0 = 0$

Using:
\[
\sin z = z - \frac{z^3}{3!} + \cdots
\]

We get:
\[
f(z) = \frac{z}{z(1 - z^2/3! + \cdots)} = \frac{1}{1 - z^2/6 + \cdots}
\]

Hence $\lim_{z \to z_0} = 1$ exist, so $z=0$ is a regular point.

\paragraph{Example}
\[
f(z) = \frac{\sin z}{2z^3 + z^4}
\]

Using expansions:
\[
\sin z = z - \frac{z^3}{6} + \cdots
\]

We get:
\[
f(z) = \frac{1}{z^2} \cdot (\text{holomorphic nonzero function})
\]

So pole of order 2 at $z=0$.

